\documentclass[11pt, a4paper]{article}

% --- Pacotes Fundamentais ---
\usepackage[utf8]{inputenc}
\usepackage[T1]{fontenc}
\usepackage[brazil]{babel}
\usepackage{geometry}
\usepackage{titlesec}
\usepackage{fancyhdr}
\usepackage{xcolor}
\usepackage{booktabs}
\usepackage{tabularx}
\usepackage{colortbl}
\usepackage{float}
\usepackage{longtable}
\usepackage{xltabular}
\usepackage{parskip}
\usepackage{lmodern}
\usepackage{enumitem}
\usepackage{amssymb}
\usepackage{tikz}
\usepackage{hyperref}
\usepackage{lastpage}

% --- Definição de Cores (idêntica ao guia anterior) ---
\definecolor{primaryColor}{RGB}{0, 51, 102}    % Azul Marinho
\definecolor{secondaryText}{RGB}{80, 80, 80}   % Cinza Escuro
\definecolor{accentColor}{RGB}{178, 34, 34}    % Vermelho Tijolo
\definecolor{lightGray}{RGB}{245, 245, 245}
\definecolor{tableHeader}{RGB}{230, 230, 250}
\definecolor{successGreen}{RGB}{34, 139, 34}
\definecolor{warningOrange}{RGB}{255, 140, 0}

% --- Configuração de Hyperlinks ---
\hypersetup{
    colorlinks=true,
    linkcolor=primaryColor,
    urlcolor=accentColor,
    pdftitle={Briefing Cotidiano - Pre-assinatura Sebraetec},
    pdfauthor={Info Dental}
}

% --- Layout de Página ---
\geometry{left=2cm, right=2cm, top=3.8cm, bottom=2cm, headsep=1.2cm, headheight=50pt}

% --- Cabeçalho e Rodapé ---
\pagestyle{fancy}
\fancyhf{}

\fancyhead[C]{%
    \begin{tabularx}{\textwidth}{@{} X r @{}}
        \textcolor{primaryColor}{\textbf{\footnotesize BRIEFING DE ALINHAMENTO --- COTIDIANO ACELERADORA DE STARTUPS S/A}} & \scriptsize \textit{Info Dental} \\
        \textcolor{secondaryText}{\scriptsize Pré-requisitos e perguntas antes da assinatura --- Proposta Sebraetec GO202604136905} & {\scriptsize \textbf{Maio 2026}}
    \end{tabularx}%
}

\fancyfoot[L]{\scriptsize \textcolor{secondaryText}{Info Dental \,|\, CNPJ 48.818.166/0001-46 \,|\, Documento interno de preparação}}
\fancyfoot[R]{\normalsize \textbf{\thepage/\pageref{LastPage}}}

\renewcommand{\headrulewidth}{1.5pt}
\renewcommand{\headrule}{%
    \vspace{0.1cm}
    \hbox to\headwidth{\color{primaryColor}\leaders\hrule height \headrulewidth\hfill}%
}

\renewcommand{\footrulewidth}{1.5pt}
\renewcommand{\footrule}{%
    \color{primaryColor}%
    \hbox to\headwidth{\leaders\hrule height \footrulewidth\hfill}%
    \vspace{0.2cm}
}

% --- Formatação de Títulos ---
\titleformat{\section}
{\color{primaryColor}\normalfont\large\bfseries}
{}{0em}{}[{\titlerule[0.8pt]}]
\titlespacing*{\section}{0pt}{10pt}{4pt}

\titleformat{\subsection}
{\color{primaryColor}\normalfont\normalsize\bfseries}
{}{0em}{}
\titlespacing*{\subsection}{0pt}{8pt}{2pt}

\setlist{nosep, leftmargin=1.4em}

% --- Caixa de destaque ---
\newcommand{\destaque}[2]{%
    \vspace{2pt}
    \colorbox{#1!10}{%
        \parbox{\dimexpr\textwidth-2\fboxsep}{%
            \small #2%
        }%
    }
    \vspace{4pt}
}

\begin{document}

%==============================================================
\section{Contexto e Objetivo Final}
%==============================================================

A Info Dental abriu chamado no SEBRAE Goiás para se capacitar e atender ao requisito alternativo da \textbf{FINEP --- Mais Inovação Brasil --- Rodada 2 --- Saúde / Empresas / Linha V --- Dispositivos Médicos} (submissão até 31/08/2026, fluxo contínuo). Em resposta, o SEBRAE-GO encaminhou a proposta \textbf{Sebraetec \texttt{GO202604136905} --- Desenvolvimento de Negócios Inovadores --- Tração e escala no mercado}, no valor total de \textbf{R\$ 20.000,00} (subsídio SEBRAE 70\% = R\$ 14.000,00; contrapartida Info Dental = R\$ 6.000,00), a ser executada pela \textbf{COTIDIANO ACELERADORA DE STARTUPS S/A} (Brasília-DF) em quatro etapas entre 04/05/2026 e 03/08/2026.

\destaque{accentColor}{%
    \textcolor{accentColor}{\textbf{Objetivo final inequívoco:}} a participação no Sebraetec é \emph{meio}, não fim. O fim é (1) gerar \textbf{documentação comprobatória} de que a Info Dental foi \textbf{selecionada como cliente em programa de inovação do SEBRAE com apoio financeiro} (requisito alternativo da FINEP ao faturamento de R\$ 1 milhão); e (2) apoiar a submissão do projeto \textbf{OralRisk.AI} ao edital FINEP até 31/08/2026, ainda dentro da janela útil do Sebraetec.%
}

%==============================================================
\section{Análise Sintética da Proposta Recebida}
%==============================================================

\renewcommand{\arraystretch}{1.25}
\begin{table}[H]
    \centering
    \small
    \begin{tabularx}{\textwidth}{@{} l X r @{}}
        \toprule
        \rowcolor{tableHeader} \textbf{Etapa} & \textbf{Descrição} & \textbf{Valor} \\
        \midrule
        \textbf{1. Alinhamento da Proposta} & Reunião de abertura, nivelamento de escopo, cronograma resumido, definição de responsáveis. \textit{04/05 a 18/05/2026}. & R\$ 5.000 \\
        \rowcolor{lightGray}
        \textbf{2. Diagnóstico e Plano de Ação} & Diagnóstico em quatro dimensões: modelo de negócios; produtos e operações; marketing e vendas; desenvolvimento organizacional. \textit{18/05 a 01/06/2026}. & R\$ 5.000 \\
        \textbf{3. Desenvolvimento do Negócio Inovador} & Reuniões semanais de mentoria, orientação, validação e suporte técnico/operacional nas quatro dimensões. \textit{01/06 a 29/06/2026}. & R\$ 5.000 \\
        \rowcolor{lightGray}
        \textbf{4. Acesso a Crédito e Investimentos} & Qualificação para participar de pelo menos um demoday, conexão com investidores e clientes. \textit{29/06 a 03/08/2026}. & R\$ 5.000 \\
        \midrule
        \multicolumn{2}{r}{\textbf{Total}} & \textbf{R\$ 20.000} \\
        \bottomrule
    \end{tabularx}
\end{table}

\destaque{warningOrange}{%
    \textcolor{warningOrange}{\textbf{Janela crítica:}} a Etapa 4 encerra em 03/08/2026 e a submissão FINEP encerra em 31/08/2026 --- restam \textbf{28 dias úteis} para finalizar comprovantes, preencher o FAP, gravar o vídeo de 10 min e protocolar. Qualquer atraso operacional do prestador inviabiliza a janela.%
}

%==============================================================
\section{Pré-requisitos Críticos --- a Resolver ANTES da Assinatura}
%==============================================================

Os itens abaixo são \textbf{condições materiais} sem as quais a assinatura do contrato Sebraetec é uma despesa de R\$ 6.000,00 sem retorno em relação ao objetivo FINEP.

\renewcommand{\arraystretch}{1.3}
{\small
\begin{xltabular}{\textwidth}{@{} l X c @{}}
    \toprule
    \rowcolor{tableHeader} \textbf{Pré-requisito} & \textbf{Por que é crítico} & \textbf{Status} \\
    \midrule
    \endfirsthead
    \toprule
    \rowcolor{tableHeader} \textbf{Pré-requisito} & \textbf{Por que é crítico} & \textbf{Status} \\
    \midrule
    \endhead
    \bottomrule
    \endlastfoot
    \textbf{PR1.} Documento comprobatório nominal &
    A FINEP exige prova formal de que a Info Dental foi \emph{selecionada como cliente em programa de inovação do SEBRAE com apoio financeiro}. O comprobatório precisa nomear a empresa, citar o subsídio (R\$ 14.000), e identificar o programa como Sebraetec --- Negócios Inovadores. &
    \textcolor{accentColor}{\textbf{Exigir}} \\
    \rowcolor{lightGray}
    \textbf{PR2.} Aderência do escopo ao OralRisk.AI &
    As 4 dimensões da proposta (negócio / produto / marketing / organizacional) precisam acomodar um SaMD em desenvolvimento --- não um \textit{pivot} forçado. O plano de ação da Etapa 2 deve ser \textbf{compatível} com a tese de produto OralRisk.AI e com o cronograma do edital Linha V. &
    \textcolor{accentColor}{\textbf{Negociar}} \\
    \textbf{PR3.} Confidencialidade (NDA) prévia &
    A Etapa 2 exige acesso ao modelo de negócio, dados financeiros, roadmap tecnológico (fine-tuning Sabiá-3, RAG, corpus clínico de $\sim$1.210 documentos). O corpus clínico do Prof.\ Cláudio M.\ Pereira é \textbf{ativo proprietário irrepetível} e diferencial competitivo central. &
    \textcolor{accentColor}{\textbf{Assinar antes}} \\
    \rowcolor{lightGray}
    \textbf{PR4.} Propriedade intelectual preservada &
    Qualquer entregável do Sebraetec (planos, diagnósticos, pitch deck, materiais de demoday) deve ser de \textbf{titularidade da Info Dental}. A Cotidiano é prestadora de serviço, não cossócia da PI. &
    \textcolor{accentColor}{\textbf{Cláusula}} \\
    \textbf{PR5.} Cronograma garantido em contrato &
    As datas 04/05 $\to$ 03/08/2026 precisam ser \textbf{cláusula contratual} (não estimativa). Atraso de qualquer etapa inviabiliza a submissão FINEP de 31/08/2026. &
    \textcolor{accentColor}{\textbf{Vincular}} \\
    \rowcolor{lightGray}
    \textbf{PR6.} Mentor nomeado com experiência em healthtech / regulado &
    Diagnóstico genérico de aceleradora de startups não serve para um SaMD ANVISA. O mentor precisa ter histórico em saúde digital, dispositivos médicos ou healthtech regulada. &
    \textcolor{accentColor}{\textbf{Exigir}} \\
    \textbf{PR7.} AFE ANVISA --- requisito independente &
    AFE ANVISA é \textbf{eliminatório próprio da Linha V}, \emph{não} resolvido pelo Sebraetec. A Info Dental precisa verificar se já tem AFE compatível com software para saúde, ou iniciar solicitação \textbf{em paralelo} (30--90 dias úteis). &
    \textcolor{warningOrange}{\textbf{Paralelo}} \\
    \rowcolor{lightGray}
    \textbf{PR8.} Cláusula de rescisão sem penalidade &
    Se ao final da Etapa 1 ficar evidente que a Cotidiano não consegue gerar o comprobatório FINEP nos moldes exigidos, a Info Dental deve poder rescindir sem multa, pagando apenas a etapa entregue. &
    \textcolor{accentColor}{\textbf{Cláusula}} \\
\end{xltabular}
}

%==============================================================
\section{Perguntas Pertinentes para a Cotidiano --- Reunião de Alinhamento}
%==============================================================

\subsection{4.1 Sobre o documento comprobatório FINEP (perguntas P1--P4)}

\begin{enumerate}[label=\textcolor{primaryColor}{\textbf{P\arabic*.}}, leftmargin=2.2em]
    \item Ao final do contrato, qual é o \textbf{documento formal} emitido pelo SEBRAE-GO (e/ou pela Cotidiano) que comprova a participação da Info Dental como \emph{cliente selecionada em programa de inovação com apoio financeiro}? Existe modelo ou exemplo de outras empresas atendidas?
    \item Esse comprobatório nomeia explicitamente: (a) razão social e CNPJ da Info Dental; (b) o programa (Sebraetec --- Negócios Inovadores --- Tração); (c) o valor do subsídio SEBRAE (R\$ 14.000,00); (d) o período de atendimento; (e) a assinatura do SEBRAE-GO?
    \item Em quanto tempo após a Etapa 4 (encerramento em 03/08/2026) o comprobatório fica disponível? É possível obter \textbf{declarações parciais} ao final de cada etapa, para uso antecipado em outros editais?
    \item A Cotidiano já atendeu empresas que usaram esse comprobatório para submissão à FINEP, MCTI, BNDES ou similares? Em caso positivo, é possível conversar com \textbf{uma referência} ou ver o documento (anonimizado) que foi aceito?
\end{enumerate}

\subsection{4.2 Sobre o escopo das 4 etapas e adequação ao OralRisk.AI (P5--P10)}

\begin{enumerate}[label=\textcolor{primaryColor}{\textbf{P\arabic*.}}, leftmargin=2.2em, start=5]
    \item Considerando que o projeto-âncora da Info Dental é o \textbf{OralRisk.AI} --- um Software como Dispositivo Médico (SaMD) sob a RDC ANVISA 657/2022 --- o diagnóstico da Etapa 2 nas quatro dimensões (modelo de negócios / produto / marketing-vendas / desenvolvimento organizacional) será \textbf{adaptado ao contexto de healthtech regulada}, ou segue um \textit{template} genérico de aceleradora?
    \item Quem será o(s) consultor(es) nominalmente alocado(s)? Pode-se enviar previamente o currículo / LinkedIn e histórico em projetos de software médico, healthtech ou edital FINEP?
    \item As mentorias semanais da Etapa 3 podem incluir orientação direta ao \textbf{preenchimento do FAP da FINEP}, à composição do orçamento de R\$ 5--15 MM da Linha V e à montagem do dossiê de submissão? Ou esse trabalho está fora do escopo Sebraetec?
    \item A Etapa 4 fala em \textit{demoday} e conexão com investidores. Para uma empresa em vias de submeter projeto FINEP, é possível \textbf{redirecionar} essa etapa para: preparação do vídeo obrigatório de 10 minutos do edital, pitch para a banca FINEP, ou pré-validação documental do FAP? Se sim, fica registrado no plano de ação da Etapa 2?
    \item O plano de ação produzido na Etapa 2 será \textbf{validado em conjunto} com a Info Dental antes da execução? A Info Dental tem poder de \textbf{veto} sobre direcionamentos que conflitem com o roadmap do OralRisk.AI?
    \item Há \textbf{entregáveis físicos / digitais nominados} ao final de cada etapa (relatório de diagnóstico, plano de ação, deck de pitch, ata de mentoria etc.)? Pode-se anexar essa lista ao contrato?
\end{enumerate}

\subsection{4.3 Sobre confidencialidade e propriedade intelectual (P11--P14)}

\begin{enumerate}[label=\textcolor{primaryColor}{\textbf{P\arabic*.}}, leftmargin=2.2em, start=11]
    \item A Cotidiano aceita assinar \textbf{NDA específico} antes da Etapa 2, cobrindo: roadmap do OralRisk.AI, corpus clínico, arquitetura técnica (Sabiá-3 + LoRA + RAG), dados financeiros, lista de clientes e parceiros estratégicos (incluindo PUC Goiás)?
    \item Como a Cotidiano trata \textbf{conflito de interesse} caso esteja acelerando, no mesmo período, outra startup do setor de saúde digital ou odontologia?
    \item Os entregáveis (diagnóstico, plano de ação, materiais de pitch, apresentações para investidores) são de \textbf{titularidade exclusiva da Info Dental}, com cessão integral de direitos? Há cláusula contratual nesse sentido?
    \item A Cotidiano se compromete a \textbf{não usar} cases / materiais / nome da Info Dental em marketing ou outros pitchs sem autorização prévia por escrito?
\end{enumerate}

\subsection{4.4 Sobre cronograma, atrasos e janela FINEP (P15--P17)}

\begin{enumerate}[label=\textcolor{primaryColor}{\textbf{P\arabic*.}}, leftmargin=2.2em, start=15]
    \item As datas 04/05/2026 $\to$ 03/08/2026 são \textbf{vinculantes em contrato}? Há cláusula de penalidade ou redução de honorários em caso de atraso de etapa atribuível à Cotidiano?
    \item Se a Etapa 1 (alinhamento) demonstrar inadequação entre o escopo Sebraetec e o objetivo FINEP, há possibilidade de \textbf{rescisão sem multa}, pagando apenas a Etapa 1 entregue?
    \item As reuniões serão \textbf{presenciais em Goiânia} ou remotas? Qual a frequência mínima das mentorias semanais da Etapa 3? Como ficam registradas (atas, gravações)?
\end{enumerate}

\subsection{4.5 Sobre experiência da Cotidiano com healthtech e FINEP (P18--P20)}

\begin{enumerate}[label=\textcolor{primaryColor}{\textbf{P\arabic*.}}, leftmargin=2.2em, start=18]
    \item Quantas startups de \textbf{healthtech / software médico / SaMD} a Cotidiano já acelerou? Quantas dessas submeteram editais FINEP, MCTI, BNDES ou similares? Quantas foram aprovadas?
    \item A Cotidiano conhece o \textbf{edital Mais Inovação Brasil --- Linha V}? Já apoiou clientes na submissão de propostas para esse edital ou similar (saúde / dispositivos médicos / subvenção econômica)?
    \item A Cotidiano tem rede de contatos relevantes para o OralRisk.AI: empresas com \textbf{AFE ANVISA vigente} (potenciais parceiras / coexecutoras), investidores em healthtech regulado, especialistas regulatórios ANVISA, NITs de outras universidades?
\end{enumerate}

\subsection{4.6 Sobre apoio adicional à documentação financeira para o FAP (P21--P24)}

A documentação financeira da Info Dental (\textbf{Balanço Patrimonial, DRE, comprovação de Patrimônio Líquido positivo} do último exercício) é \textbf{critério eliminatório} de habilitação na FINEP --- mas \emph{não} é, formalmente, escopo do Sebraetec. Como healthtech \textbf{pré-breakeven} (RO negativo natural à fase de P\&D) e com \textbf{time majoritariamente PJ} (não CLT) --- perfil típico de startup tech brasileira ---, a Info Dental enfrenta três desafios estruturantes: garantir PL positivo, estruturar contrapartida sem queimar caixa, e \emph{enquadrar} o time PJ de forma aceitável pela FINEP. Faz sentido sondar a Cotidiano, que trabalha rotineiramente com manobras societárias e estruturações financeiras de startups em estágio similar.

\begin{enumerate}[label=\textcolor{primaryColor}{\textbf{P\arabic*.}}, leftmargin=2.2em, start=21]
    \item Embora a contabilidade não seja escopo formal do Sebraetec, a Cotidiano poderia, dentro das mentorias da Etapa 3, \textbf{orientar} a Info Dental na organização do Balanço e DRE --- incluindo, se necessário, manobras pré-submissão como \textbf{aumento de capital social} (para destravar PL positivo, se atualmente negativo por prejuízos acumulados) e estruturação de \textbf{contrapartida via pessoal próprio}, considerando que o time é majoritariamente PJ?
    \item A Cotidiano dispõe, em sua rede, de \textbf{contador ou consultor contábil} com experiência específica em empresas de base tecnológica \textbf{pré-breakeven} que apresentam balanços a editais de fomento (FINEP, BNDES, EMBRAPII), familiarizado com aumento de capital, conversão de SAFE / mútuos conversíveis em equity e Balanço Patrimonial sob regime contábil completo (não Simples)? Pode indicar profissionais?
    \item Como apoio adicional, a Cotidiano pode \textbf{revisar} --- durante a Etapa 2 (Diagnóstico) --- a coerência entre: (a) orçamento previsto do projeto FINEP (R\$ 5--15 MM, contrapartida 5--10\%); (b) estrutura de contrapartida (folha + infraestrutura dedicada + aporte em conta exclusiva); e (c) indicadores financeiros eliminatórios da Info Dental (PL positivo, endividamento oneroso $\leq$30\% Ativo Total se RO negativo, contrapartida $\leq$50\% Ativo Total), sinalizando inconsistências bloqueantes antes da submissão?
    \item A Cotidiano tem experiência em orientar startups quanto ao \textbf{enquadramento de colaboradores PJ vs CLT} para fins de contrapartida FINEP? Especificamente: (a) \emph{conversão seletiva} de 1--3 colaboradores-chave (tech lead, cientista de dados sênior, gerente de projeto) para CLT imediatamente antes da submissão; (b) \emph{modelagem reforçada} de contratos PJ (longa duração, vinculação explícita ao projeto, dedicação documentada, cláusulas de exclusividade parcial) com chance maior de serem aceitos como pessoal próprio; (c) eventual consulta formal escrita à FINEP (\texttt{cp\_mib\_sbv\_saude\_2@finep.gov.br}) para validar a interpretação \emph{antes} da submissão?
\end{enumerate}

%==============================================================
\section{Cláusulas Recomendadas no Termo de Adesão / Contrato}
%==============================================================

\begin{enumerate}[label=\textcolor{primaryColor}{\textbf{C\arabic*.}}, leftmargin=2.2em]
    \item \textbf{Objetivo declarado:} consignar expressamente que a finalidade contratada é \emph{capacitar a Info Dental para submissão de projeto de inovação à FINEP até 31/08/2026}, gerando o comprobatório de cliente Sebraetec --- e que essa finalidade é elemento essencial do contrato.
    \item \textbf{Confidencialidade (NDA):} cláusula de sigilo com prazo mínimo de 5 anos, cobrindo informações técnicas, financeiras, comerciais, corpus clínico e roadmap.
    \item \textbf{Propriedade intelectual:} titularidade integral dos entregáveis para a Info Dental; cessão automática e gratuita de quaisquer direitos autorais sobre relatórios, planos, decks e materiais produzidos.
    \item \textbf{Vinculação de cronograma:} datas de início e fim de cada etapa como obrigação contratual; descontos proporcionais em caso de atraso atribuível ao prestador.
    \item \textbf{Rescisão sem multa após Etapa 1:} faculdade unilateral da Info Dental de rescindir após a Etapa 1, sem ônus além do valor da etapa entregue, caso seja constatada inadequação ao objetivo declarado.
    \item \textbf{Mentor nomeado:} identificação do(s) consultor(es) responsável(is) no contrato, com vedação a substituição sem aprovação prévia da Info Dental.
    \item \textbf{Entregáveis nominados:} anexo contratual listando o produto físico/digital de cada etapa, com critérios de aceitação.
    \item \textbf{Não-concorrência limitada:} declaração da Cotidiano de inexistência de conflito de interesse com outras aceleradas no mesmo nicho durante a execução.
\end{enumerate}

%==============================================================
\section{Guia Operacional Info Dental --- Itens Fora do Escopo Sebraetec}
%==============================================================

O Sebraetec resolve \textbf{apenas um} dos requisitos da Linha V (o comprobatório de empresa selecionada em programa de inovação). Os \textbf{cinco itens} abaixo são \emph{responsabilidade direta da Info Dental} e devem correr em \textbf{paralelo} às 4 etapas do Sebraetec (04/05 $\to$ 03/08/2026). Cada um traz pré-requisitos, custos, prazos e responsáveis explícitos. Ignorar qualquer deles é fatal para a habilitação FINEP.

%--------------------------------------------------------------
\subsection{6.1 AFE ANVISA --- Autorização de Funcionamento de Empresa}

\destaque{accentColor}{%
    \textcolor{accentColor}{\textbf{ELIMINATÓRIO Linha V.}} Ao menos uma das beneficiárias do arranjo (proponente ou coexecutora) deve possuir AFE ANVISA \textbf{vigente} no ato da submissão. Sem AFE: a proposta é eliminada na habilitação, sem recurso.%
}

\begin{itemize}[label=\textcolor{primaryColor}{$\blacktriangleright$}, leftmargin=1.6em]
    \item \textbf{O que é:} autorização da ANVISA para empresa fabricar, importar ou distribuir produtos para saúde --- incluindo \textit{software médico} (SaMD) sob a RDC 657/2022.
    \item \textbf{Como obter:} portal \texttt{gov.br/anvisa} $\to$ Serviços $\to$ \emph{Empresas} $\to$ \emph{Peticionamento eletrônico de AFE}.
    \item \textbf{Documentos necessários:}
    \begin{itemize}[label=$\circ$, leftmargin=1.4em]
        \item Contrato social atualizado, com objeto social compatível com produto para saúde / software médico.
        \item CNPJ ativo e regular; certidões negativas federal, estadual, municipal e FGTS.
        \item Comprovante de endereço da sede e do local de operação.
        \item Indicação de \textbf{Responsável Técnico} habilitado conforme RDC 16/2014 (para software médico: profissional de saúde ou engenheiro biomédico com registro de classe ativo).
        \item Pagamento da GRU (Guia de Recolhimento da União) referente à taxa ANVISA.
    \end{itemize}
    \item \textbf{Prazo realista:} 30--90 dias úteis após protocolo (varia conforme demanda da ANVISA e necessidade de cumprimento de exigências).
    \item \textbf{Custo estimado:} taxa ANVISA \emph{(GRU)} R\$ 200--R\$ 5.000 conforme porte (ME pode ter desconto ou isenção, art.\ 6º da Lei nº 11.972/2009). Consultoria regulatória especializada (recomendada): R\$ 5.000--R\$ 15.000.
    \item \textbf{Responsável Info Dental:} sócios + consultor regulatório ANVISA (contratar se não tiver internamente).
    \item \textcolor{accentColor}{\textbf{AÇÃO IMEDIATA (esta semana):}} verificar se a Info Dental \emph{já possui} AFE vigente. Se sim, conferir se o objeto social cobre \textit{software para saúde}. Se não, protocolar pedido em até \textbf{15 de junho de 2026} (margem para 75 dias até a submissão FINEP em 31/08).
    \item \textbf{Plano B:} identificar \emph{empresa parceira} já com AFE vigente para entrar no arranjo como coexecutora --- via NIT/PUC Goiás (pergunta 1 da reunião com NIT) ou via rede da Cotidiano (P20).
\end{itemize}

%--------------------------------------------------------------
\subsection{6.2 Arranjo com ICT --- NIT da PUC Goiás}

\begin{itemize}[label=\textcolor{primaryColor}{$\blacktriangleright$}, leftmargin=1.6em]
    \item \textbf{O que é:} instrumento contratual entre empresa proponente (Info Dental) e a ICT (PUC Goiás) formalizando a parceria de pesquisa, desenvolvimento e inovação --- amparo no \textbf{Marco Legal da Inovação} (Lei nº 13.243/2016) e no Decreto nº 9.283/2018.
    \item \textbf{Como obter:} agendar reunião com o NIT (Núcleo de Inovação Tecnológica) da PUC Goiás.
    \item \textbf{Contato:}
    \begin{itemize}[label=$\circ$, leftmargin=1.4em]
        \item Prof.\ Dr.\ Marcos Lajovic Carneiro --- Coordenador do NIT.
        \item WhatsApp: \textbf{(62) 98130-6931}.
        \item E-mail: \texttt{mcarneiro@pucgoias.edu.br}.
    \end{itemize}
    \item \textbf{O que negociar com o NIT:}
    \begin{itemize}[label=$\circ$, leftmargin=1.4em]
        \item Tipo de instrumento: \emph{Acordo de Parceria para PD\&I} (não convênio).
        \item Cláusula de \textbf{Serviços de Consultoria} para remunerar o Prof.\ Cláudio Maranhão Pereira (vedação a bolsa e pró-labore no edital FINEP) --- ICT emite NF para a empresa proponente.
        \item Percentual mínimo do orçamento para a ICT: $\geq$15\%.
        \item Cláusula de PI compartilhada (\% a definir; titularidade dos modelos de IA e do corpus tratada explicitamente).
        \item Cronograma de atividades exclusivas da ICT (curadoria do corpus, validação clínica IEC 62304, avaliação cega do LLM, orientação de 4 alunos IC).
        \item Cláusula de confidencialidade.
    \end{itemize}
    \item \textbf{Prazo realista:} 30--60 dias entre primeira reunião e assinatura --- aprovação interna da PUC depende da reitoria.
    \item \textbf{Custo direto:} zero (a ICT é remunerada via recursos do projeto FINEP).
    \item \textbf{Responsável Info Dental:} representante legal + Prof.\ Cláudio M.\ Pereira (já vinculado à PUC).
    \item \textcolor{accentColor}{\textbf{AÇÃO IMEDIATA:}} agendar reunião com o NIT esta semana, levando: ficha do edital (já elaborada), proposta OralRisk.AI v2 (já elaborada), CNPJ da Info Dental.
\end{itemize}

%--------------------------------------------------------------
\subsection{6.3 CEP --- Comitê de Ética em Pesquisa}

\begin{itemize}[label=\textcolor{primaryColor}{$\blacktriangleright$}, leftmargin=1.6em]
    \item \textbf{O que é:} aprovação ética para uso de dados clínicos de pacientes (anonimizados, conforme LGPD) na curadoria do corpus de fine-tuning do Sabiá-3 e na validação dos modelos de IA. \textbf{Condição prévia} à curadoria de casos clínicos do banco proprietário do Prof.\ Cláudio M.\ Pereira (50+ MRONJ, 108+ anticoagulados).
    \item \textbf{Onde submeter:} \textbf{Plataforma Brasil} --- \texttt{plataformabrasil.saude.gov.br} --- direcionando ao CEP da PUC Goiás.
    \item \textbf{Documentos necessários:}
    \begin{itemize}[label=$\circ$, leftmargin=1.4em]
        \item Projeto de pesquisa detalhado (uso secundário de dados clínicos para treinamento de IA).
        \item Protocolo de \textbf{anonimização LGPD} conforme ABNT NBR ISO/IEC 29101.
        \item TCLE (Termo de Consentimento Livre e Esclarecido) \emph{ou} pedido de \textbf{dispensa de TCLE} para uso secundário (banco já existente, anonimização irreversível).
        \item Currículo Lattes do pesquisador principal + co-pesquisadores.
        \item Cronograma e local de execução (Clínica Escola de Odontologia da PUC Goiás).
    \end{itemize}
    \item \textbf{Prazo realista:} 60--90 dias (análise inicial em $\sim$30 dias + eventuais pendências). \textbf{Curadoria do corpus não pode iniciar antes da aprovação.}
    \item \textbf{Custo:} zero.
    \item \textbf{Responsável Info Dental:} formalmente nenhum --- a submissão é feita pelo Prof.\ Cláudio M.\ Pereira (pesquisador principal) com apoio do NIT/PUC. Info Dental acompanha como interessada.
    \item \textcolor{accentColor}{\textbf{AÇÃO IMEDIATA (Maio--Junho 2026):}} alinhar com o pesquisador a submissão à Plataforma Brasil em até 30 dias --- aprovação ideal antes do Mês 4 do cronograma do projeto.
\end{itemize}

%--------------------------------------------------------------
\subsection{6.4 Cadastro FINEP --- Plataforma de Captação}

\begin{itemize}[label=\textcolor{primaryColor}{$\blacktriangleright$}, leftmargin=1.6em]
    \item \textbf{O que é:} cadastro obrigatório da empresa proponente na plataforma da FINEP, condição prévia ao preenchimento do FAP (Formulário de Apresentação de Propostas).
    \item \textbf{Onde:} \texttt{cadastro.finep.gov.br} --- preenchimento online + envio digital de documentos.
    \item \textbf{Documentos necessários:}
    \begin{itemize}[label=$\circ$, leftmargin=1.4em]
        \item CNPJ ativo, contrato social atualizado, última alteração contratual.
        \item Dados do(s) sócio(s): CPF, RG, comprovante de endereço.
        \item Dados do representante legal (procuração se aplicável).
        \item Dados bancários da conta da empresa (banco / agência / conta) --- para liberação futura dos recursos.
        \item Balanço Patrimonial e DRE do último exercício (mesmos documentos da seção 6.5).
    \end{itemize}
    \item \textbf{Prazo realista:} 5--10 dias úteis para validação pela FINEP, podendo ser maior se houver exigências.
    \item \textbf{Custo:} zero.
    \item \textbf{Responsável Info Dental:} representante legal.
    \item \textcolor{accentColor}{\textbf{AÇÃO IMEDIATA (Mês 2 do Sebraetec, $\sim$junho 2026):}} iniciar cadastro \emph{antes} de finalizar o FAP, pois a plataforma de submissão (\texttt{financiamento.finep.gov.br}) só libera o formulário após cadastro aprovado.
\end{itemize}

%--------------------------------------------------------------
\subsection{6.5 Documentação Financeira --- Capacidade de Contrapartida}

\destaque{primaryColor}{%
    \textcolor{primaryColor}{\textbf{Dois critérios FINEP independentes --- não confundir:}} \\[2pt]
    \textbf{(A) Entrada / elegibilidade da empresa:} faturamento $\geq$ R\$ 1 MM \emph{ou} comprobatório de seleção em programa de inovação SEBRAE com apoio financeiro. \textcolor{successGreen}{$\to$ \emph{resolvido pelo Sebraetec.}} \\
    \textbf{(B) Capacidade financeira (habilitação):} PL positivo + regras de contrapartida vs.\ Ativo Total / RO. \textcolor{accentColor}{$\to$ \emph{NÃO resolvido pelo Sebraetec --- responsabilidade da Info Dental.}}\\[2pt]
    Faturamento baixo é problema do critério A. PL negativo é problema do critério B. \textbf{São barreiras diferentes, com soluções diferentes.}%
}

\destaque{accentColor}{%
    \textcolor{accentColor}{\textbf{ELIMINATÓRIO Linha V.}} A análise de capacidade financeira da empresa proponente é critério eliminatório na habilitação. Os documentos abaixo devem estar prontos e \textbf{coerentes} antes do preenchimento final do FAP.%
}

\begin{itemize}[label=\textcolor{primaryColor}{$\blacktriangleright$}, leftmargin=1.6em]
    \item \textbf{O que é:} prova documental de que a Info Dental tem condição financeira de aportar a contrapartida exigida (5\% ME / 10\% PE sobre o valor total) e operar o projeto com saúde financeira.
    \item \textbf{Documentos exigidos:}
    \begin{itemize}[label=$\circ$, leftmargin=1.4em]
        \item \textbf{Balanço Patrimonial} do último exercício social encerrado, assinado por contador habilitado (CRC).
        \item \textbf{DRE} (Demonstração de Resultado do Exercício) do último exercício, assinada por contador.
        \item Comprovação de \textbf{Patrimônio Líquido positivo} (deriva do Balanço; eliminatório).
        \item Declaração de faturamento ou Receita Operacional Bruta (ROB) --- enquadra a empresa no porte (ME / EPP / PE) e define o \% mínimo de contrapartida.
        \item Certidões negativas: federal, estadual, municipal, FGTS, trabalhista.
    \end{itemize}
    \item \textbf{Critérios financeiros FINEP (todos eliminatórios):}
    \begin{itemize}[label=$\circ$, leftmargin=1.4em]
        \item Patrimônio Líquido \textbf{positivo}.
        \item Se Resultado Operacional \textbf{positivo}: contrapartida $\leq$ 20\% do RO.
        \item Se Resultado Operacional \textbf{negativo}: endividamento oneroso $\leq$ 30\% do Ativo Total \emph{e} contrapartida $\leq$ 50\% do Ativo Total.
    \end{itemize}
    \item \textbf{Onde obter:} contador atual da Info Dental. Se a empresa é optante pelo Simples Nacional e só mantém escrituração simplificada, \textbf{contratar contador para emitir Balanço sob regime contábil completo} (formato exigido pela FINEP).
    \item \textbf{Prazo realista:} 15--30 dias se o Balanço completo ainda não estiver pronto; 5--10 dias se já existe.
    \item \textbf{Custo estimado:} R\$ 1.500--R\$ 5.000 para emissão de Balanço completo + DRE auditados (varia conforme contador e complexidade); R\$ 0 se já produzidos rotineiramente.
    \item \textbf{Responsável Info Dental:} sócios + contador.
    \item \textbf{Apoio Cotidiano (a confirmar com P21--P23):} eventual orientação dentro das mentorias da Etapa 3 + indicação de contador parceiro da rede da aceleradora.
    \item \textcolor{accentColor}{\textbf{AÇÃO IMEDIATA (esta semana):}} pedir ao contador atual o \textbf{último Balanço e DRE} emitidos. Se a Info Dental \textbf{não tem Balanço completo} dos últimos 12 meses, encomendar imediatamente.
\end{itemize}

\vspace{4pt}
\textbf{\textcolor{primaryColor}{Caminho específico para healthtech pré-breakeven (caso Info Dental):}}

\vspace{2pt}
RO negativo é \textbf{esperado} para healthtech em fase de P\&D e \emph{não} elimina por si só --- existe regra alternativa via Ativo Total. O risco real está no \textbf{PL negativo} (prejuízos acumulados $>$ capital social $+$ reservas) e na \textbf{ausência de caixa para contrapartida em dinheiro}. Ambos têm contorno padrão:

\begin{itemize}[label=\textcolor{accentColor}{$\blacktriangleright$}, leftmargin=1.6em]
    \item \textbf{Passo 1 --- diagnóstico PL.} Olhar o último Balanço: se PL atual já é positivo (capital social $+$ aportes $-$ prejuízos acumulados $>$ 0), o critério está atendido --- prosseguir. Se negativo, executar Passo 2.
    \item \textbf{Passo 2 --- aumento de capital social} (manobra padrão pré-edital). Sócios e/ou anjos aportam capital novo via instrumento de aumento de capital social, registrado em alteração contratual na Junta Comercial. Novo Balanço pós-aumento reflete PL positivo. Custo $\sim$R\$ 500--R\$ 1.500 (Junta + contador); prazo $\sim$30 dias.
    \item \textbf{Passo 3 --- conversão de SAFE / mútuo conversível} em equity, se houver. Operação societária análoga ao aumento de capital, com mesmo efeito sobre o PL.
    \item \textbf{Passo 4 --- estruturar contrapartida via pessoal próprio (atenção ao perfil CLT vs PJ).} Em proposta de R\$ 5 MM (mínimo Linha V), contrapartida ME (5\%) = R\$ 250 k em 36 meses $\approx$ R\$ 7 k/mês. As regras de enquadramento:
    \begin{itemize}[label=$\circ$, leftmargin=1.4em]
        \item \textbf{CLT alocado em P\&D:} aceito sem ambiguidade como contrapartida. \textcolor{successGreen}{$\checkmark$ Caminho seguro.}
        \item \textbf{PJ típico (NF mensal, prestação de serviços):} entra no orçamento como \emph{Serviços de Terceiros} --- financiável pelo recurso FINEP, \emph{mas não substitui contrapartida}. \textcolor{warningOrange}{$\triangle$ Não conta para os 5\%.}
        \item \textbf{PJ "robusto"} (contrato longo, vinculação explícita ao projeto, dedicação documentada, exclusividade parcial): zona cinzenta --- pode ser aceito como pessoal próprio em alguns editais, com risco de reclassificação para subcontratação (que é vedada pelo art.\ 4.4 do Regulamento). \textcolor{accentColor}{$!$ Validar com consulta escrita à FINEP.}
    \end{itemize}
    \item \textbf{Passo 4b --- como o time da Info Dental é majoritariamente PJ}, as três rotas práticas são:
    \begin{itemize}[label=$\circ$, leftmargin=1.4em]
        \item \textbf{(a) Conversão seletiva CLT} de 1--3 chaves (tech lead, cientista de dados sênior, gerente de projeto). Encargos CLT $\sim$70\% sobre o salário bruto. Garantia jurídica máxima --- recomendado para os perfis nucleares.
        \item \textbf{(b) Modelagem PJ reforçada} + consulta formal escrita à FINEP --- risco controlado se houver resposta documentada antes da submissão.
        \item \textbf{(c) Contrapartida via infraestrutura dedicada} (servidores GPU, licenças de software, equipamentos): pode somar R\$ 100--200 k em 36 meses, complementando o pessoal e reduzindo a pressão sobre a folha.
        \item \textbf{Estrutura recomendada Info Dental:} 2 CLT seniores ($\sim$R\$ 15 k/mês cada $\Rightarrow$ $\sim$R\$ 540 k em folha bruta em 36 meses) + restante do time como PJ na rubrica Serviços de Terceiros (paga pelo FINEP) + infra GPU complementar. Cobre folgadamente os R\$ 250 k de contrapartida.
    \end{itemize}
    \item \textbf{Passo 5 --- documentação de suporte.} Anexar ao FAP, além de Balanço e DRE: cap table atualizado, contratos de investimento anjo (se houver), projeção de fluxo de caixa do projeto, \textbf{minutas dos contratos CLT e PJ} dos colaboradores alocados ao projeto, carta de comprometimento de aporte futuro de sócios (opcional, mas reforça a tese de viabilidade).
\end{itemize}

\destaque{successGreen}{%
    \textcolor{successGreen}{\textbf{Bottom line para Info Dental:}} pré-breakeven \emph{não} é impeditivo --- inúmeras healthtechs aprovam FINEP nessa condição. Os três pontos que matam são: (1) submeter com PL negativo \emph{sem} ter feito aumento de capital antes; (2) estruturar contrapartida em caixa que a empresa não tem, em vez de pessoal próprio + infraestrutura; (3) tentar usar a folha PJ inteira como contrapartida \emph{sem} converter ao menos 1--3 chaves para CLT \emph{ou} obter resposta escrita da FINEP validando o modelo PJ reforçado. Todos têm solução conhecida e barata, desde que iniciadas \emph{agora} (margem de $\sim$60 dias para aumento de capital + conversões CLT seletivas + novo Balanço).%
}

%--------------------------------------------------------------
\subsection{6.6 Resumo: cronograma consolidado dos itens paralelos}

\renewcommand{\arraystretch}{1.3}
\begin{table}[H]
    \centering
    \small
    \begin{tabularx}{\textwidth}{@{} l X c c @{}}
        \toprule
        \rowcolor{tableHeader} \textbf{Item} & \textbf{Ação Imediata} & \textbf{Prazo} & \textbf{Custo} \\
        \midrule
        AFE ANVISA & Verificar AFE atual; se ausente, protocolar até 15/06 & 30--90 dias & R\$ 5--20k \\
        \rowcolor{lightGray}
        Arranjo ICT (NIT) & Reunião com Prof.\ Marcos Lajovic esta semana & 30--60 dias & R\$ 0 \\
        CEP (Plataforma Brasil) & Submissão até final de junho & 60--90 dias & R\$ 0 \\
        \rowcolor{lightGray}
        Cadastro FINEP & Iniciar em junho/2026 & 5--10 dias & R\$ 0 \\
        Documentação financeira & Solicitar Balanço + DRE ao contador esta semana & 15--30 dias & R\$ 1,5--5k \\
        Aumento de capital (se PL$<$0) & Diagnóstico do PL atual; se negativo, aporte + alteração contratual & $\sim$30 dias & R\$ 0,5--1,5k \\
        \rowcolor{lightGray}
        Conversão seletiva PJ$\to$CLT & Identificar 1--3 chaves; preparar contratos CLT e migração & 15--30 dias & Encargos $\sim$70\% s/ salário \\
        Consulta escrita à FINEP (opcional) & Pergunta formal sobre enquadramento PJ; guardar resposta & 15--45 dias & R\$ 0 \\
        \bottomrule
    \end{tabularx}
\end{table}

%==============================================================
\section{Decisão: assinar, renegociar ou recusar}
%==============================================================

\renewcommand{\arraystretch}{1.3}
\begin{table}[H]
    \centering
    \small
    \begin{tabularx}{\textwidth}{@{} l X @{}}
        \toprule
        \rowcolor{tableHeader} \textbf{Cenário} & \textbf{Ação recomendada} \\
        \midrule
        \textcolor{successGreen}{\textbf{Assinar}} & Cotidiano confirma: (a) comprobatório FINEP no formato exigido; (b) mentor com experiência em healthtech; (c) escopo adaptável ao OralRisk.AI; (d) NDA; (e) PI da Info Dental; (f) cronograma vinculante; (g) rescisão sem multa após Etapa 1. \\
        \rowcolor{lightGray}
        \textcolor{warningOrange}{\textbf{Renegociar}} & Cotidiano oferece escopo genérico mas aceita inclusão das cláusulas C1--C8 e aceita redirecionar a Etapa 4 (\textit{demoday}) para apoio à submissão FINEP. \\
        \textcolor{accentColor}{\textbf{Recusar / pedir outro credenciado}} & Cotidiano não consegue (ou não quer) emitir comprobatório nominado para FINEP; ou exige escopo genérico de aceleradora sem aderência a healthtech; ou recusa NDA / PI / rescisão. Voltar ao SEBRAE-GO e solicitar novo credenciado. \\
        \bottomrule
    \end{tabularx}
\end{table}

\destaque{primaryColor}{%
    \textcolor{primaryColor}{\textbf{Resumo de uma linha:}} a Info Dental deve assinar somente após a Cotidiano confirmar, por escrito, que o comprobatório emitido pelo SEBRAE-GO atende ao requisito da FINEP de \emph{empresa selecionada como cliente em programa de inovação do SEBRAE com apoio financeiro}, e após o contrato incorporar as cláusulas C1--C8 acima.%
}

\end{document}
